\UseRawInputEncoding
\documentclass[a4paper, 14pt]{extarticle}

% Поля
%--------------------------------------
\usepackage{geometry}
\geometry{a4paper,tmargin=2cm,bmargin=2cm,lmargin=3cm,rmargin=1cm}
%--------------------------------------


%Russian-specific packages
%--------------------------------------
\usepackage[T2A]{fontenc}
\usepackage[utf8]{inputenc} 
\usepackage[english, main=russian]{babel}
%--------------------------------------

\usepackage{textcomp}

% Красная строка
%--------------------------------------
\usepackage{indentfirst}               
%--------------------------------------             


%Graphics
%--------------------------------------
\usepackage{graphicx}
\graphicspath{ {./images/} }
\usepackage{wrapfig}
%--------------------------------------

% Полуторный интервал
%--------------------------------------
\linespread{1.3}                    
%--------------------------------------

%Выравнивание и переносы
%--------------------------------------
% Избавляемся от переполнений
\sloppy
% Запрещаем разрыв страницы после первой строки абзаца
\clubpenalty=10000
% Запрещаем разрыв страницы после последней строки абзаца
\widowpenalty=10000
%--------------------------------------

%Списки
\usepackage{enumitem}

%Подписи
\usepackage{caption} 

%Гиперссылки
\usepackage{hyperref}

\hypersetup {
	unicode=true
}

%Рисунки
%--------------------------------------
\DeclareCaptionLabelSeparator*{emdash}{~--- }
\captionsetup[figure]{labelsep=emdash,font=onehalfspacing,position=bottom}
%--------------------------------------

\usepackage{tempora}
\usepackage{amsmath}
\usepackage{color}
\usepackage{listings}
\lstset{
  belowcaptionskip=1\baselineskip,
  breaklines=true,
  frame=L,
  xleftmargin=\parindent,
  language=Python,
  showstringspaces=false,
  basicstyle=\footnotesize\ttfamily,
  keywordstyle=\bfseries\color{blue},
  commentstyle=\itshape\color{purple},
  identifierstyle=\color{black},
  stringstyle=\color{red},
  extendedchars=\true,
}

%--------------------------------------
%			НАЧАЛО ДОКУМЕНТА
%--------------------------------------

\begin{document}

%--------------------------------------
%			ТИТУЛЬНЫЙ ЛИСТ
%--------------------------------------
\begin{titlepage}
\thispagestyle{empty}
\newpage


%Шапка титульного листа
%--------------------------------------
\vspace*{-60pt}
\hspace{-65pt}
\begin{minipage}{0.3\textwidth}
\hspace*{-20pt}\centering

\includegraphics[width=\textwidth]{emblem}
\end{minipage}
\begin{minipage}{0.67\textwidth}\small \textbf{
\vspace*{-0.7ex}
\hspace*{-6pt}\centerline{Министерство науки и высшего образования Российской Федерации}
\vspace*{-0.7ex}
\centerline{Федеральное государственное бюджетное образовательное учреждение }
\vspace*{-0.7ex}
\centerline{высшего образования}
\vspace*{-0.7ex}
\centerline{<<Московский государственный технический университет}
\vspace*{-0.7ex}
\centerline{имени Н.Э. Баумана}
\vspace*{-0.7ex}
\centerline{(национальный исследовательский университет)>>}
\vspace*{-0.7ex}
\centerline{(МГТУ им. Н.Э. Баумана)}}
\end{minipage}
%--------------------------------------

%Полосы
%--------------------------------------
\vspace{-25pt}
\hspace{-35pt}\rule{\textwidth}{2.3pt}

\vspace*{-20.3pt}
\hspace{-35pt}\rule{\textwidth}{0.4pt}
%--------------------------------------

\vspace{1.5ex}
\hspace{-35pt} \noindent \small ФАКУЛЬТЕТ\hspace{80pt} <<Информатика и системы управления>>

\vspace*{-16pt}
\hspace{47pt}\rule{0.83\textwidth}{0.4pt}

\vspace{0.5ex}
\hspace{-35pt} \noindent \small КАФЕДРА\hspace{50pt} <<Теоретическая информатика и компьютерные технологии>>

\vspace*{-16pt}
\hspace{30pt}\rule{0.866\textwidth}{0.4pt}
  
\vspace{11em}

\begin{center}
\Large {\bf Лабораторная работа № 9} \\ 
\large {\bf<<Реализация сингулярного разложения матрицы (SVD)>>} \\ 
\large {\bf по курсу <<Численные методы линейной алгебры>>} \\ 
\end{center}\normalsize

\vspace{8em}


\begin{flushright}
  {Студент группы ИУ9-71Б Лысенко В. В.\hspace*{15pt} \\
  \vspace{2ex}
  Преподаватель Посевин Д. П.\hspace*{15pt}}
\end{flushright}

\bigskip

\vfill
 

\begin{center}
\textsl{Москва 2024}
\end{center}
\end{titlepage}
%--------------------------------------
%		КОНЕЦ ТИТУЛЬНОГО ЛИСТА
%--------------------------------------

\renewcommand{\ttdefault}{pcr}

\setlength{\tabcolsep}{3pt}
\newpage
\setcounter{page}{2}

\section{Цель работы}\label{Sect::task}
\par
\renewcommand{\labelenumii}{\arabic{enumi}.\arabic{enumii}.}
Самостоятельно изучить сингулярное разложения матрицы или SVD разложение.

\section{Задачи}\label{Sect::task}
\par
\renewcommand{\labelenumii}{\arabic{enumi}.\arabic{enumii}.}
\begin{enumerate} 
    \item Реализовать SVD разложение матрицы.
    \item Проверить корректность реализации прямым перемножением матриц.
\end{enumerate}

\pagebreak
\section{Практическая реализация}

\par
Код программы:
\begin{lstlisting}
using LinearAlgebra


function Mysvd(A)
    m, n = size(A)
    
    mat = A' * A
    
    eigvals, V = eigen(mat)
    
    idx = sortperm(eigvals, rev=true)
    eigvals = eigvals[idx]
    V = V[:, idx]
    
    S = sqrt.(eigvals)
    
    lambda = zeros(m, n)
    for i in 1:min(m, n)
        lambda[i, i] = S[i]
    end
    
    U = zeros(m, m)
    for i in 1:size(V, 2)
        U[:, i] = A * V[:, i] / S[i]
    end
    
    return U, lambda, V
end

function calculate_error(A, A_approx)
    return norm(A - A_approx)
end


n = 5
A = rand(n, n) * 100
println("Матрица A:")
for row in eachrow(A)
    for x in row
        print(x, "  ")
    end
    println()
end
println()

U, lambda, V = Mysvd(A)

println("Матрица U:")
for row in eachrow(U)
    for x in row
        print(x, "  ")
    end
    println()
end
println()




println("Матрица lambda:")
for row in eachrow(lambda)
    for x in row
        print(x, "  ")
    end
    println()
end
println()

println("Матрица V:")
for row in eachrow(V)
    for x in row
        print(x, "  ")
    end
    println()
end
println()



A_approx = U * lambda * V'
println("Востановленная матрица A:")
for row in eachrow(A_approx)
    for x in row
        print(x, "  ")
    end
    println()
end
println()

error = calculate_error(A, A_approx)
println("Погрешность разложения:")
println(error)
\end{lstlisting}

\section{Результаты}
\par
Результат выполнения программы:
\begin{center}
\includegraphics[scale=0.55]{g1}
\end{center}
\par


\section{Выводы}
В ходе выполнения лаборатороной работы был самостоятельно изучен алгоритм сингулярного разложения матрицы или SVD разложения.
\end{document}
